\begin{table}[!htbp]
\raggedright
\caption{Momentary profile of the four challenge-skill channels.}
\label{tab:sup-channelprofiles}
\footnotesize
\begingroup
\setlength{\tabcolsep}{4pt}
\renewcommand{\arraystretch}{1.08}
\noindent
\begin{tabular*}{\textwidth}{@{\extracolsep{\fill}}lcccccccc@{}}
\toprule
Channel & $n$ & \% & Pain & Interference & Attention & Threat & Flow exp. & Activation \\
\midrule
Flow & 1,314 & 22.0 & 4.90 & 4.15 & 2.99 & 2.34 & 5.79 & 4.33 \\
Relaxation & 2,214 & 37.0 & 4.29 & 3.52 & 2.69 & 2.09 & 4.64 & 2.47 \\
Anxiety & 516 & 8.6 & 4.83 & 4.28 & 3.11 & 2.23 & 4.44 & 4.44 \\
Apathy & 1,932 & 32.3 & 4.49 & 3.94 & 3.00 & 2.31 & 3.33 & 2.46 \\
\bottomrule
\end{tabular*}
\par\vspace{3pt}\noindent\parbox{\textwidth}{\footnotesize\textit{Note.} Unadjusted momentary means on the raw 1 to 7 scale under the strict absolute rule (both condition items at least 5). Activation is the self-reported physical activation item, which is far higher in the two high-challenge channels and is the reason every model reported in the text is also estimated with activation as a covariate.}
\endgroup
\end{table}
